\documentclass[10pt,a4paper]{amsart}
\usepackage{amsmath,amssymb,amsthm,graphicx}
% Retained minimal preamble for the Stacks Project chapter more-algebra.tex (More on Algebra).
% The chapter file pinned-original-source.tex carries only \input{preamble} on its first line, so no
% preamble text lives inside the chapter itself.  The Stacks Project uses one single shared preamble for
% every chapter; the snapshot stages it identically as sources/stacks-chow/preamble.tex and
% sources/stacks-etale-cohomology/preamble.tex.  The block below is copied mechanically, line for line,
% from sources/stacks-chow/preamble.tex lines 14-18 (the Xy-pic setup) and 154-217 (the native theorem
% declarations together with the mathematical macros).  Nothing was added, removed or redefined; the
% wrapper class and the amsmath/amssymb/amsthm/graphicx base packages are supplied by the assembler.
% Chapter: more-algebra.tex -- More on Algebra.
% For commutative diagrams we use Xy-pic
\usepackage[all]{xy}

% We use 2cell for 2-commutative diagrams.
\xyoption{2cell}

% Theorem environments.
%
\theoremstyle{plain}
\newtheorem{theorem}[subsection]{Theorem}
\newtheorem{proposition}[subsection]{Proposition}
\newtheorem{lemma}[subsection]{Lemma}

\theoremstyle{definition}
\newtheorem{definition}[subsection]{Definition}
\newtheorem{example}[subsection]{Example}
\newtheorem{exercise}[subsection]{Exercise}
\newtheorem{situation}[subsection]{Situation}

\theoremstyle{remark}
\newtheorem{remark}[subsection]{Remark}
\newtheorem{remarks}[subsection]{Remarks}

\numberwithin{equation}{subsection}

% Macros
%
\def\lim{\mathop{\mathrm{lim}}\nolimits}
\def\colim{\mathop{\mathrm{colim}}\nolimits}
\def\Spec{\mathop{\mathrm{Spec}}}
\def\Hom{\mathop{\mathrm{Hom}}\nolimits}
\def\Ext{\mathop{\mathrm{Ext}}\nolimits}
\def\SheafHom{\mathop{\mathcal{H}\!\mathit{om}}\nolimits}
\def\SheafExt{\mathop{\mathcal{E}\!\mathit{xt}}\nolimits}
\def\Sch{\mathit{Sch}}
\def\Mor{\mathop{\mathrm{Mor}}\nolimits}
\def\Ob{\mathop{\mathrm{Ob}}\nolimits}
\def\Sh{\mathop{\mathit{Sh}}\nolimits}
\def\NL{\mathop{N\!L}\nolimits}
\def\CH{\mathop{\mathrm{CH}}\nolimits}
\def\proetale{{pro\text{-}\acute{e}tale}}
\def\etale{{\acute{e}tale}}
\def\QCoh{\mathit{QCoh}}
\def\Ker{\mathop{\mathrm{Ker}}}
\def\Im{\mathop{\mathrm{Im}}}
\def\Coker{\mathop{\mathrm{Coker}}}
\def\Coim{\mathop{\mathrm{Coim}}}

% Boxtimes
%
\DeclareMathSymbol{\boxtimes}{\mathbin}{AMSa}{"02}

%
% Macros for moduli stacks/spaces
%
\def\QCohstack{\mathcal{QC}\!\mathit{oh}}
\def\Cohstack{\mathcal{C}\!\mathit{oh}}
\def\Spacesstack{\mathcal{S}\!\mathit{paces}}
\def\Quotfunctor{\mathrm{Quot}}
\def\Hilbfunctor{\mathrm{Hilb}}
\def\Curvesstack{\mathcal{C}\!\mathit{urves}}
\def\Polarizedstack{\mathcal{P}\!\mathit{olarized}}
\def\Complexesstack{\mathcal{C}\!\mathit{omplexes}}
% \Pic is the operator that assigns to X its picard group, usage \Pic(X)
% \Picardstack_{X/B} denotes the Picard stack of X over B
% \Picardfunctor_{X/B} denotes the Picard functor of X over B
\def\Pic{\mathop{\mathrm{Pic}}\nolimits}
\def\Picardstack{\mathcal{P}\!\mathit{ic}}
\def\Picardfunctor{\mathrm{Pic}}
\def\Deformationcategory{\mathcal{D}\!\mathit{ef}}
\title{Every finite set of elements of a bounded above complex of free modules lies in a bounded subcomplex of finite free modules}
\author{The Stacks Project Authors}
\date{stacks-project:more-algebra.tex:local-snapshot-20261012; GNU FDL 1.2 or later}
\begin{document}
\maketitle
\setcounter{subsection}{1}
\noindent\textit{Editorial note.} All delivered mathematical text is the author's own, copied from the pinned Stacks Project chapter file more-algebra.tex; no statement or proof step was rewritten, reordered or removed. The counted statement is the labeled lemma environment at source lines 38552-38559 (canonical tag 0FNR) and the counted proof is the proof environment at source lines 38561-38578, which immediately follows it and is its own proof.\par\medskip
\noindent\textit{Editorial note.} The lead-in sentence at source lines 38548-38550, announcing that a simple lemma is needed before the dual and invertible objects of D(R) can be analysed, is outside the counted spans because it belongs to the surrounding discussion. The lemma that follows this one, at source lines 38580-38591, is a different statement and is not counted.\par\medskip
\noindent\textit{Editorial note.} The chapter file pinned-original-source.tex carries only an input of the shared preamble on its first line, so no preamble text lives inside the chapter. The retained file sources/stacks-more-algebra/retained-preamble.tex is a mechanical, line-for-line copy of the project's single shared preamble as staged in the snapshot (Xy-pic setup plus the native theorem declarations and mathematical macros). Nothing was added, removed or redefined, and the excerpt serves only as preamble, never as a span coordinate.\par\medskip
\noindent\textit{Editorial note.} The source URL is the source-level Stacks Project URL, as required when more than one card is produced from one source; the per-claim tag is carried by claim\_key and by this note. The counted statement and proof refer to no labelled item of the chapter and cite no bibliography entry, so this card contributes no external reference stub and needs no bibliography.\par\medskip
\noindent\textit{Allowed standard prerequisites (editorial).}\begin{itemize}
\item An element of a direct sum of free modules on one generator each involves only finitely many summands, so the elements of a finite list lie in the span of a finite subset of a chosen basis. used at proof lines 38567-38571, where the finite subsets J' of J and K' of K are chosen. Produces the finite rank pieces out of which the bounded subcomplex is built.
\item Descending induction on the lowest degree of the given elements, with the zero complex as the base case when the complex vanishes above that degree. used at proof lines 38562-38566 and again at proof lines 38572-38575. Reduces the statement for a finite list of elements to the statement for the elements of strictly higher degree after the degree a piece has been chosen.
\item A subcomplex may be specified by giving its terms in degrees above a and one more term in degree a, provided the differential of that last term lands in the next term. used at proof lines 38571-38577, where the differential of the basis elements b\_j is required to land in the span of the chosen c\_k. Justifies taking G equal to H in degrees above a and to the span of the b\_j in degree a, so that G is a subcomplex with finite free terms and bounded support.
\end{itemize}\medskip
\begin{lemma}
\label{lemma-complex-bounded-above-free-colim-bounded-finite-free}
Let $R$ be a ring. Let $F^\bullet$ be a bounded above complex of
free $R$-modules. Given pairs $(n_i, f_i)$, $i = 1, \ldots, N$
with $n_i \in \mathbf{Z}$ and $f_i \in F^{n_i}$ there exists
a subcomplex $G^\bullet \subset F^\bullet$ containing
all $f_i$ which is bounded and consists of finite free $R$-modules.
\end{lemma}

\begin{proof}
By descending induction on $a = \min(n_i; i = 1, \ldots, N)$.
If $F^n = 0$ for $n \geq a$, then the result is true with
$G^\bullet$ equal to the zero complex. In general, after renumbering
we may assume there exists an $1 \leq r \leq N$ such that
$n_1 = \ldots = n_r = a$ and $n_i > a$ for $i > r$.
Choose a basis $b_j, j \in J$ for $F^a$. We can choose a finite subset
$J' \subset J$ such that $f_i \in \bigoplus_{j \in J'} Rb_j$
for $i = 1, \ldots, r$. Choose a basis $c_k, k \in K$ for $F^{a + 1}$.
We can choose a finite subset $K' \subset K$ such that
$\text{d}_F^a(b_j) \in \bigoplus_{k \in K'} Rc_k$ for
$j \in J'$. Then we can apply the induction
hypothesis to find a subcomplex $H^\bullet \subset F^\bullet$
containing $c_k \in F^{a + 1}$ for $k \in K'$ and
$f_i \in F^{n_i}$ for $i > r$. Take $G^\bullet$ equal to
$H^\bullet$ in degrees $> a$ and equal to $\bigoplus_{j \in J'} Rb_j$
in degree $a$.
\end{proof}

\end{document}
